\documentclass[10pt]{article}

\usepackage[margin=0.85in]{geometry}
\usepackage{amsmath,amssymb,booktabs,enumitem,microtype}
\usepackage[hidelinks]{hyperref}

\newcommand{\clip}{\operatorname{clip}}
\newcommand{\CE}{\operatorname{CE}}
\newcommand{\Normal}{\mathcal{N}}
\newcommand{\Dpub}{D_{\mathrm{pub}}}
\newcommand{\Dupd}{D_{\mathrm{upd}}}
\newcommand{\Dprobe}{D_{\mathrm{probe}}}

\title{Plan for a Counterfactual Clipping-Advantage Dataset}
\author{WRN-16-1 / CIFAR-10 DP-SGD experiment}
\date{23 September 2026}

\begin{document}
\maketitle

\begin{abstract}
This document specifies a one-step counterfactual dataset for learning a
per-example adaptive clipping policy. Starting from each of 5,100 stored
DP-SGD checkpoints, two reproducible Poisson batches are drawn from a declared
public subset of CIFAR-10. Two examples per batch are selected as anchors. A
baseline update with clipping norm one is compared with updates in which only
one anchor uses either a low or a high clipping value. The label is the paired
change in cross-entropy on a disjoint public probe set. The resulting dataset
contains exactly 40,800 rows.
\end{abstract}

\section{Inputs and privacy premise}

Let $D$ be the 50,000-example CIFAR-10 training set. There are $L=10$
independent runs and $T=510$ stored post-update checkpoints per run:
\[
  \Theta=\{\theta_{s,\ell}:s\in\{1,\ldots,T\},
  \ell\in\{0,\ldots,L-1\}\}.
\]
The original configuration is WRN-16-1, $q=1/3$, clipping norm $C=1$,
noise multiplier $\sigma=4.888246618$, learning rate $\eta=4.0$, no momentum,
and no weight decay. The source collection is specified in
\texttt{DATA-COLLECTION.md}.

A fixed 10\% subset
\[
  \Dpub\subset D,\qquad |\Dpub|=5{,}000,
\]
is declared public side information before construction. Its indices and
selection seed are recorded. Only examples, labels, gradients, and histories
belonging to $\Dpub$ may be used. The remaining 45,000 examples and their
saved per-example features are out of scope.

The public subset is partitioned once into disjoint sets
\[
  \Dpub=\Dupd\mathbin{\dot\cup}\Dprobe,
  \qquad |\Dupd|=n_{\mathrm{upd}}=4{,}000,
  \qquad |\Dprobe|=n_{\mathrm{probe}}=1{,}000.
\]
$\Dupd$ supplies counterfactual update batches and anchors. $\Dprobe$ is used
only for the probe loss and never belongs to a counterfactual update batch.
The split is class-stratified, and its seed and exact indices are recorded.

The privacy premise is that the checkpoint transcript $\Theta$ is an
$(\varepsilon,\delta)$-DP release with respect to the non-public records. Under
this premise, the construction is post-processing of $\Theta$, public data,
and data-independent randomness, and incurs no additional privacy loss for the
non-public records. This conclusion does not hold if the builder accesses
stored gradients or histories of non-public examples.

\section{Sampling design}

For every checkpoint $(s,\ell)$, draw two independent Poisson batches:
\[
  \mathcal{B}_{s\ell r}
  =\{j\in\Dupd:Z_{s\ell rj}=1\},
  \qquad Z_{s\ell rj}\sim\operatorname{Bernoulli}(q),
  \qquad q=\frac13,\quad r\in\{1,2\}.
\]
Therefore
\[
  \mathbb{E}|\mathcal{B}_{s\ell r}|=q|\Dupd|
  =4{,}000/3.
\]
The seeds vary with $(s,\ell,r)$ and are derived deterministically from a
master seed. The two batches may overlap.

For each batch select $b=2$ distinct anchors uniformly without replacement:
\[
  I_{s\ell r}=\{i_{s\ell r1},i_{s\ell r2}\}
  \subset\mathcal{B}_{s\ell r}.
\]
The action set is
\[
  \mathcal{C}=\{C_{\mathrm{low}},C_{\mathrm{high}}\}=\{0.25,0.75\}.
\]
The baseline and both actions obey the global cap
$C_{\mathrm{cap}}=1$.

\section{Paired counterfactual}

For a per-example gradient $g_j(\theta)$, define
\[
  \clip(g_j,C)=g_j\min\left\{1,\frac{C}{\|g_j\|_2}\right\}.
\]
Let $N=50{,}000$ denote the deployment population size. The public batch
estimates the background clipped mean,
\[
  \overline G_{s\ell r}^{(1)}
  =\frac{1}{q n_{\mathrm{upd}}}
    \sum_{j\in\mathcal B_{s\ell r}}\clip(g_j(\theta_{s,\ell}),1).
\]
The noise retains the original deployment scale:
\[
  \xi_{s\ell r}\sim\Normal(0,\sigma^2C_{\mathrm{cap}}^2I),
  \qquad
  \theta_{s\ell r}^{(1)}
  =\theta_{s,\ell}-\eta\left(
      \overline G_{s\ell r}^{(1)}+\frac{\xi_{s\ell r}}{qN}
    \right).
\]
For anchor $i\in I_{s\ell r}$ and action $a\in\mathcal C$, define the
one-record deployment-scale perturbation
\[
  \Delta G_{s\ell ria}
  =\frac{\clip(g_i(\theta_{s,\ell}),a)
          -\clip(g_i(\theta_{s,\ell}),1)}{qN},
  \qquad
  \theta_{s\ell r}^{(i,a)}
  =\theta_{s\ell r}^{(1)}-\eta\Delta G_{s\ell ria}.
\]
Thus the 4,000 public update examples estimate the background direction, while
the individual action and DP noise retain their original $N=50{,}000$ scale.
The baseline and counterfactuals use identical batch membership,
augmentations, optimizer state, and Gaussian noise. Different $(s,\ell,r)$
tuples use different recorded seeds.

The independent public probe objective is
\[
  L_{\mathrm{probe}}(\theta)
  =\frac{1}{|\Dprobe|}
    \sum_{(x,y)\in\Dprobe}\CE(f_\theta(x),y),
\]
evaluated on canonical, non-augmented examples. The paired advantage is
\[
  A_{s\ell ria}
  =L_{\mathrm{probe}}\!\left(\theta_{s\ell r}^{(1)}\right)
   -L_{\mathrm{probe}}\!\left(\theta_{s\ell r}^{(i,a)}\right).
\]
Thus $A>0$ is a reward, $A<0$ is a penalty, and $A=0$ is neutral. The baseline
next state is computed once per $(s,\ell,r)$ and shared by four
counterfactuals. All temporary next states are discarded after evaluation.

A diagnostic subset additionally uses $a=1$. Its advantage must be numerically
zero under the paired-randomness protocol.

\section{Rows and dataset size}

The stored checkpoint index is $s$. A logical decision at $t=s+1$ is recorded
as
\[
  (\theta_{t-1,\ell},h_i^{1:t-1},a,A_{ita})
  =(\theta_{s,\ell},h_i^{1:s},a,A_{s\ell ria}).
\]
$h_i^{1:s}$ references the existing sequence of 4,275-dimensional projected
gradients and six exact block norms for public anchor $i$. Checkpoints and
histories are referenced, not copied into every row.

The number of primary rows is
\[
  L\,T\,(2\ \text{batches})\,(b\ \text{anchors})\,|\mathcal C|
  =10\cdot510\cdot2\cdot2\cdot2
  =\boxed{40{,}800}.
\]

Each row stores at least
\begin{center}
\begin{tabular}{ll}
\toprule
Field & Meaning \\
\midrule
\texttt{run\_id}, \texttt{checkpoint\_step} & Key for $\theta_{s,\ell}$ \\
\texttt{checkpoint\_uri}, \texttt{history\_ref} & Feature references \\
\texttt{batch\_replicate}, \texttt{anchor\_slot} & $r\in\{1,2\}$ and anchor 1 or 2 \\
\texttt{example\_id}, \texttt{candidate\_clip} & Public anchor and action $a$ \\
\texttt{baseline\_probe\_ce} & $L_{\mathrm{probe}}(\theta^{(1)})$ \\
\texttt{candidate\_probe\_ce} & $L_{\mathrm{probe}}(\theta^{(i,a)})$ \\
\texttt{advantage\_raw} & $A_{s\ell ria}$ \\
\texttt{batch\_seed}, \texttt{augmentation\_seed}, \texttt{noise\_seed} & RNG metadata \\
\texttt{public\_split\_hash} & Identity of $\Dupd$ and $\Dprobe$ \\
\bottomrule
\end{tabular}
\end{center}
All 40,800 primary rows are retained without filtering by sign or magnitude.

\section{Post-construction scaling}

The raw advantage is authoritative. After construction, define
\[
  s_A=\operatorname{median}_{s,\ell,r,i,a}|A_{s\ell ria}|,
  \qquad
  \widetilde A_{s\ell ria}
  =\frac{A_{s\ell ria}}{s_A+10^{-12}}.
\]
This single global multiplier improves numerical scale without changing signs
or action rankings. Store $A$, $\widetilde A$, and $s_A$. Do not rescale
separately by step or run.

\section{Validation and limitations}

Before full construction, verify:
\begin{enumerate}[leftmargin=*,nosep]
  \item $\Dupd\cap\Dprobe=\varnothing$ and
        $\Dupd\cup\Dprobe=\Dpub$;
  \item every anchor belongs to its recorded batch and to $\Dupd$;
  \item only declared-public sample indices are accessed;
  \item paired baseline and counterfactual RNG states agree exactly;
  \item diagnostic action $a=1$ has zero advantage to numerical tolerance;
  \item recorded seeds reproduce sampled indices and labels.
\end{enumerate}

The label is a local, one-step marginal advantage around the $C=1$ policy; it
does not identify the long-horizon return of changing many clipping decisions.
The fixed public probe loss is the stated optimization target, not an estimate
of performance on an unseen test distribution. Two actions do not identify an
optimal continuous clipping value. Any learned policy must therefore be
evaluated in fresh end-to-end DP-SGD runs.

\end{document}
